

% !TeX root = graphical.tex  
\chapter{Bayesian Networks}
\label{chap:bayes.net}
\section{Definitions}

Bayesian networks are graphical models supported
by directed acyclic graphs (DAG), which provide them with an ordered
organization (directed graphs were introduced in  \cref{def:dir.grph}). 

We first introduce some notation. Let $G=(V,E)$ be a directed acyclic graph. The parents of $s\in V$ are
vertexes $t$ such that $(t,s)\in E$, and its children are $t$'s such
that $(s,t)\in E$. The set of parents of $s$ is denoted $\pa{s}$, and
the set of its children is $\cl{s}$, with $\CV_s = \cl{s} \cup
\pa{s}$.

Similarly to trees, the vertexes of $G$ can be
partially ordered by $s\leq_G t$ if and only if
there exists a path going from $s$ to $t$. Unlike trees, however,
there can be more than one minimal element in $V$, and we
still call roots vertexes that have no parent, denoting
$$
V_0 = \defset{s\in V: \pa{s} = \emp}.
$$
We also call leaves, or terminal nodes, vertexes that have no children. Unless otherwise
specified, we assume that all graphs are connected.

 Bayesian networks over $G$ are defined as follows. We use the same notation as with Markov random fields to represent the set of configurations $\CF(V)$ that contains collections $x = (x_s, s\in V)$ with $x_s\in F_s$.
\begin{definition}
\label{def:bn}
A random variable $X$ with values in $\CF(V)$ is a Bayesian network over a DAG $G=(V,E)$ if and only
if its distribution can be written in the form
\begin{equation}
\label{eq:bn}
P_X(x) = \prod_{s\in V_0} p_s(\pe xs) \prod_{s\in V\setm V_0} p_s(\pe x{\pa{s}}, \pe x s)
\end{equation}
where $p_s$ is, for all $s\in V$, a probability distribution with
respect to $\pe xs$.
\end{definition}
Using the convention that conditional distributions given the empty set
are just absolute distributions, we can rewrite \cref{eq:bn} as
\begin{equation}
\label{eq:bn.2}
P_X(x) = \prod_{s\in V} p_s(\pe 
x{\pa{s}}, \pe x s).
\end{equation}
One can verify that $\sum_{x\in \Om} P^X(x) =1$. Indeed, when summing over $x$, we can
start summing over all $\pe xs$ with $\cl{s}= \emp$ (the leaves). Such
$\pe xs$'s only appear in the corresponding $p_s$'s, which disappear since they sum
to 1. What remains is the sum of the product over $V$ minus the
leaves, and the argument can be iterated until the remaining sum is 1
(alternatively, work by induction on $|V|$). This fact is also a
consequence of  \cref{prop:rest} below, applied with $A=\emp$.



\section{Conditional independence graph}
\subsection{Moral graph}
Bayesian networks have a conditional independence structure which is
not exactly given by $G$, but can be deduced from it. Indeed, fixing $S\sub V$, we can see, when computing the probability of $\pe XS=\pe xs$
given $\pe X{S^c} = \pe x{S^c}$, which is
\[
\myP(\pe XS = \pe x S \mid \pe X{S^c} = \pe x{S^c} ) = \frac{1}{Z(\pe x{S^c})} \prod_{s\in V} p_s(\pe x{\pa{s}}, \pe xs),
\]
 that the only variables $\pe xt, t\not \in S$ that
can be factorized in the normalizing constant are those that are
neither parent nor children of vertexes in $S$, and do not share a
child with a vertex in $S$ (i.e., they intervene in no
$p_s(\pe x {\pa{s}}, \pe xs)$ that  involve elements of $S$). This suggests the following definition.
\begin{definition}
\label{def:g.sharp}
Let $G$ be a directed acyclic graph. We denote
$G^\sharp=(V,E^\sharp)$ the undirected graph on $V$ such that
$\defset{s,t}\in E^\sharp$ if one of the following conditions is satisfied
\begin{itemize}
\item Either $(s,t)\in E$ or $(t,s)\in E$.
\item There exists $u\in V$ such that $(s,u)\in E$ and $(t,u)\in
E$.
\end{itemize}
\end{definition}
$G^\sharp$ is sometimes called the {\em moral graph} of $G$ (because
it forces parents to marry~!). A path in $G^\sharp$ can be visualized as a path in $G^\flat$ (the undirected graph associated with $G$) which is
allowed to jump between parents of the same vertex even if they were
not connected originally.

The previous discussion implies:
\begin{proposition}
\label{prop:g.sharp}
Let $X$ be a Bayesian network on $G$. We have
$$
\cind{S}{T}{U}_{G^\sharp}\Ria \cind{\pe XS}{\pe XT}{\pe XU},
$$
i.e., $X$ is $G^\sharp$-Markov.
\end{proposition}



This proposition can be refined by noticing that the joint
distribution of $\pe XS$, $\pe XT$ and $\pe XU$ can be deduced from a Bayesian
network on a graph restricted to the ancestors of $S\cup T\cup
U$. \Cref{def:subg} for restricted graphs extends without change to directed graphs, and we repeat it below for convenience.
\begin{definition}
\label{def:rest.grph}
Let $G=(V,E)$ be a graph (directed or undirected), and $A\sub V$. The
restricted graph $G_A = (A, E_A)$ is such that the elements of $E_A$
are the edges $(s,t)$ (or $\defset{s,t}$) in $E$ such that both $s$
and $t$ belong to $A$.
\end{definition}
Moreover, for a directed acyclic graph $G$ and $s\in V$, we define
the set of ancestors of $s$ by
\begin{equation}
\label{eq:anc}
\CA_s = \defset{t\in V, t\leq_G s}
\end{equation}
for the partial order on $V$ induced by $G$. 

If $S\sub V$, we denote $\CA_S = \bigcup_{s\in S} \CA_s$. Note that,
by definition, $S\sub \CA_S$. The following proposition is true.
\begin{proposition}
\label{prop:rest}
Let $X$ be a Bayesian network on $G=(V,E)$ with distribution given by
\cref{eq:bn.2}. Let $S\sub V$ and $A = \CA_S$. Then the distribution
of $\pe XA$ is a Bayesian network over $G_A$ given by
\begin{equation}
\label{eq:bn.rest}
\myP(\pe XA = \pe xA) = \prod_{s\in A} p_s(\pe 
x{\pa{s}}, \pe xs).
\end{equation}
\end{proposition}
There is no ambiguity in the notation $\pa{s}$, since the
parents of $s\in A$ are the same in $G_A$ as in $G$. 
\begin{proof}
One needs to show that
\[
\prod_{s\in A} p_s(\pe x{\pa{s}}, \pe xs) = \sum_{x_{A^c}} \prod_{s\in V} p_s(\pe x{\pa{s}}, \pe xs).
\]
This can be done by induction on the cardinality of $V$. Assume that
the result is true for graphs of size $n$, and let $|V| = n+1$ (the
result is obvious for graphs of size 1).

If $A=V$, there is nothing to prove, so assume that $A^c$ is not
empty. Then $A^c$ must contain a leaf in $G$, since otherwise, $A$
would contain all leaves and their ancestors which would imply that
$A=V$. 

If $s\in A^c$ is a leaf in $G$, one can remove the variable $\pe xs$ from
the sum, since it only appear in $p_s$ and transition probabilities
sum to one. But one can now apply the induction assumption to the
restriction of $G$ to $V\setm \{s\}$. 
\end{proof}

Given  \cref{prop:rest},  \cref{prop:g.sharp} can
therefore be refined as follows.
\begin{proposition}
\label{prop:g.sharp.2}
Let $X$ be a Bayesian network on $G$. We have
$$
\cind{S}{T}{U}_{(G_{\CA_{S\cup T\cup U}})^\sharp}\Ria \cind{\pe XS}{\pe XT}{\pe XU}.
$$
\end{proposition}

 \Cref{prop:rest} is also used in the proof of the following
proposition.
\begin{proposition}
\label{prop:bayes.net.cond}
Let $G = (V,E)$ be a directed acyclic graph, and $X$ be a Bayesian
network over $G$. Then, for all $s\in S$
\[
\myP(\pe Xs = \pe xs \mid \pe X{\mathcal A_s\setm \{s\}} = \pe x{\mathcal A_s\setm \{s\}}) = \myP(\pe Xs = \pe xs \mid \pe X{\pa{s}} = \pe x{s^{-}}) = p_s(\pe x{\pa{s}}, \pe xs).
\]
\end{proposition}
\begin{proof} 
By  \cref{prop:rest}, we can without loss of generality
assume that $V = \mathcal A_s$. Then
\begin{eqnarray*}
\myP(\pe Xs = \pe xs \mid \pe X{\mathcal A_s\setm \{s\}} = \pe x{\mathcal A_s\setm
  \{s\}}) &\propto & \myP(\pe X{\mathcal A_s} = \pe x{\mathcal A_s})\\
&=& p_s(\pe x{\pa{s}}, \pe xs) Z(\pe x{\mathcal A_s\setm
  \{s\}})
\end{eqnarray*}
where
\[
Z(\pe x{\mathcal A_s\setm
  \{s\}}) = \prod_{t\in \mathcal A_s\setm\{s\}} p_t(\pe x{\pa{t}}, \pe xt)
\]
disappears when the conditional probability is normalized.
\end{proof}

\subsection{Reduction to d-separation}
We now want to reformulate 
\cref{prop:g.sharp.2} in terms of the unoriented graph $G^\flat$ and
specific features in $G$ called v-junctions, that we now define.
\begin{definition}
\label{def:v.j}
Let $G=(V,E)$ be a directed graph. A v-junction is a triple of
distinct vertexes,
$(s,t,u)\in V\times V \times V$ such that $\defset{s,u}\sub \pa{t}$
(i.e., $s$ and $u$ are parents of $t$). 
   
We will say that a path $(s_1, \ldots, s_N)$ in $G^\flat$ passes at $s=s_k$ with a
v-junction if $(s_{k-1}, s_k, s_{k+1})$ is a v-junction in $G$.
\end{definition}

We have the lemma:
\begin{lemma}
\label{lem:d.sep}
Two vertexes $s$ and $t$ in $G$ are
separated by a set $U$ in $(G_{\CA_{\{s,t\}\cup U}})^\sharp$ if and
only if any path between $s$ and $t$ in $G^\flat$ must either
\begin{itemize}
\item[(1)] Pass at a vertex in $U$ without a v-junction.
\item[(2)] Pass in  $V \setm \CA_{\{s,t\}\cup U}$ at a v-junction.
\end{itemize}
\end{lemma}
\begin{proof}
\

{\noindent \it Step 1.} We first note that the v-junction clause is
redundant in (2). It can be removed without affecting the
condition. Indeed, if a path in $G^\flat$ passes in $V \setm
\CA_{\{s,t\}\cup U}$ one can follow this path downward (i.e.,
following the orientation in $G$) until a v-junction is met. This has
to happen before reaching the extremities of the path, since $u$ would
be an ancestor of $s$ or $t$ otherwise.  We can therefore work with the weaker
condition (that we will denote (2)')
in the rest of proof.

\bigskip
{\noindent \em Step 2.} 
Assume that $U$ separates $s$ and $t$ in $(G_{\CA_{\{s,t\}\cup
U}})^\sharp$. Take a path $\ga$ between $s$ and $t$ in $G^\flat$. We need to
show that the path satisfies (1) or (2)'. So assume that (2)' is false
(otherwise we are done) so that $\gamma$ is included in
$\CA_{\{s,t\}\cup U}$. We can modify $\ga$ by
removing all the central nodes in v-junctions and still keep a valid path in $(G_{\CA_{\{s,t\}\cup
U}})^\sharp$ (since parents are connected in the moral graph). The remaining path must intersect $U$ by assumption, and
this cannot be at a v-junction in $\ga$ since we have removed them. So
(1) is true.

\bigskip
\noindent {\em Step 3.}
Conversely, assume that (1) or (2) is true for any path in
$G^\flat$. Consider a path $\ga$ in $(G_{\CA_{\{s,t\}\cup U}})^\sharp$
between $s$ and $t$. Any edge in $\ga$ that is not in $G^\flat$ must
involve parents of a common child in $\CA_{\{s,t\}\cup
  U}$. Insert this child between the parents every time this occurs,
resulting in a v-junction added to $\gamma$. Since the added
vertexes are still in $\CA_{\{s,t\}\cup
  U}$, the new path still has no intersection with $V\setm \CA_{\{s,t\}\cup
  U}$ and must therefore satisfy (1). So there must be an intersection
with $U$ without a v-junction, and since the new additions are all
at v-junctions, the intersection must have been originally in $\ga$,
which therefore passes in $U$. This shows that $U$ separates $s$ and
$t$ in $(G_{\CA_{\{s,t\}\cup U}})^\sharp$.
\end{proof}

Condition (2) can be further restricted to provide the notion of
$d$-separation.
\begin{definition}
\label{def:d.sep}
One says that two vertexes $s$ and $t$ in $G$ are
$d$-separated by a set $U$ if and
only if any path between $s$ and $t$ in $G^\flat$ must either
\begin{itemize}
\item[(D1)] Pass at a vertex in $U$ without a v-junction.
\item[(D2)] Pass in  $V \setm \CA_{U}$ with a v-junction.
\end{itemize}
\end{definition}

Then we have:
\begin{theorem}
\label{th:d.sep}
Two vertexes $s$ and $t$ in $G$ are
separated by a set $U$ in $(G_{\CA_{\{s,t\}\cup U}})^\sharp$ if and
only if they are $d$-separated by $U$.
\end{theorem}
\begin{proof}
It suffices to show that if condition ((D1) or (D2)) holds for any path between
$s$ and $t$ in $G^\flat$, then so does  ((1) or (2)).  So take a path between
$s$ and $t$: if (D1) is true for this path, the
conclusion is obvious, since (D1) and (1) are the same. So assume that
(D1) (and therefore (1)) is false and that (D2) is true. Let $u$ be a
vertex in $V\setm \mathcal A_U$ at which $\ga$ passes with a v-junction.

Assume that (2) is false. Then $u$ must be an
ancestor of either $s$ or $t$. Say it is an ancestor of $s$: there is
a path in $G$ going from $u$ to $s$  without passing by $U$ (otherwise
$u$ would be an ancestor of $U$); one can replace the portion
of the old path between $s$ and $u$ by this new one, which does not
pass by $u$ with a v-junction anymore. So the new path still does
not satisfy (D1) and must satisfy (D2). Keep on removing all
intersections with ancestors of $s$ and $t$ that have v-junctions to
finally obtain a path that satisfies neither (D1) or (D2) and a
contradiction to the
fact that $s$ and $t$ are $d$-separated by $U$.
\end{proof}

\subsection{Chain-graph representation}
The $d$-separability property involves both unoriented and oriented
edges. It is in fact a property of the hybrid graph in which the
orientation is removed from the edges that are not involved in a
v-junction, and retained otherwise. Such graphs are particular
instances of chain graphs.

\begin{definition}
\label{def:ch.grph}
A chain graph $G = (V,E, \tilde E)$ is composed with a finite set $V$
of vertexes, a set $E\sub {\boldsymbol \CP}_2(V)$ of unoriented edges and a set
$\tE\sub E\ti E \setm \defset{(t,t), t\in E}$ of oriented edges with
the property that $E \cap \tE^\flat=\emp$, i.e., two vertexes cannot
be linked by both an oriented and an unoriented edge.

A path in a chain graph is a a sequence of vertexes $s_0,
\ldots, s_N$ such that for all $k\geq 1$, $s_{k-1}$ and $s_k$ form an
edge, which means that either $\defset{s_{k-1}, s_k}\in E$ or
$(s_{k-1}, s_k) \in \tE$. 

A chain graph is acyclic if it contains no loop. It is semi-acyclic if
it contains no loop containing oriented edges.
\end{definition}
 
We start with the following equivalence relation within vertexes in a semi-acyclic
chain graph.

\begin{proposition}
Let $G = (V,E, \tilde E)$ be a semi-acyclic chain graph. Define the
relation $s\,\CR\, t$ if and only if
there exists a path in the unoriented subgraph $(V,E)$ that links $s$
and $t$. Then $\CR$ is an equivalence relation.
\end{proposition}

The proposition is obvious. This relation partitions $V$ in
equivalence classes, the set
of which being denoted $V_\CR$. If $S\in V_\CR$, then any pair $s,
t$ in $S$ is related by an unoriented path, and if $S\neq S'\in
V_\CR$, no elements $s\in S$ and $t\in S'$ can be related by such a
path.

Moreover, no path in $G$ between two elements of $S\in
V_\CR$, can contain a directed edge,  since these elements must also be related by an undirected
path, and this would create a loop in $G$ containing an undirected edge. So the restriction of $G$
to $S$ is an undirected graph. 

One can define a directed graph over equivalence classes as follows. Let $G_\CR =
(V_\CR, E_\CR)$ be such that
$(S,S') \in E_\CR$ if and only if there exists $s\in S$ and $t\in
S'$ such that $(s,t)\in\tE$. The graph $G_\CR$ is acyclic: any loop in $G_\CR$
would induce a loop in $G$ containing at least one oriented edge.

We now can formally define a probability distribution on a semi-acyclic chain graph.
\begin{definition}
\label{def:acg}
Let $G = (V, E, \tE)$ be a semi-acyclic chain graph. One says that a random variable
$X$ decomposes on $G$ if and only if: $(\pe XS, S\in V_\CR)$ is
a Bayesian network on $G_\CR$ and the conditional distribution of
$\pe XS$ given $\pe X{S'}, S'\in \pa{S}$ is $G_S$-Markov, such that, for $s\in S$,
$P(\pe Xs = \pe xs\mid \pe Xt, t\in S, X_{S'}, S'\in \pa{S})$ only depends on $\pe xt$
with $\defset{s,t}\in E$ or $(t,s) \in \tilde E$.
\end{definition}


Returning to our discussion on Bayesian networks, we have the
following.  Associate to a DAG $G = (V,E)$ the chain graph
$G^\dagger = (V, E^\dagger, \tE^\dagger)$ defined by: $\defset{s,t}\in
E^\dagger$ if and only if $(s,t)$ or $(t,s)\in E$ and is not involved
in a v-junction, and $(s,t)\in \tE^\dagger$ if $(s,t)\in E$ and is
involved in a v-junction. This graph is acyclic; indeed, take any loop
in $G^\dagger$: when its edges are given their original 
orientations in $E$, the sequence cannot contain a v-junction, since
the orientation in v-junctions are kept in $G^\dagger$; the path therefore
constitutes a loop in $G$ which is a contradiction. 

All, excepted at most one, vertexes in an
equivalence class $S\in G^\dagger_\CR$ have all their parents in
$S$. Indeed, assume that two vertexes,  $s$ and $t$, in $S$ have
parents outside of
$S$. There exists an unoriented path, $s_0=s, s_1, \ldots, s_N = t$,
in $G^\dagger$ connecting them,  since they belong to the same
equivalence class. The edge
at $s$ must be oriented from $s$ to $s_1$ in $G$, since otherwise
$s_1$ would be a second parent to $s$ in $G$, creating a v-junction,
and the edge would have remained
oriented in $G^\dagger$. Similarly, the last edge in the path must be oriented from $t$ to
$s_{N-1}$ in $G$. But this implies that there exists a v-junction in
the original orientation along the path, which cannot be constituted
with only unoriented edges in $G^\dagger$. So we get a contradiction.

Thus,  random variables that decompose on $G^\dagger$ are ``Bayesian
networks'' of acyclic graphs, or trees since we know these are
equivalent. The root of each tree must have multiple (vertex) parents
in the parent tree in $G_\CR$. The following theorem states that
all Bayesian networks are equivalent to such a process.

\begin{theorem}
\label{th:acg}
Let $G = (V,E)$ be a DAG. The random variable $X$ is a Bayesian network on
$G$ if and only if it decomposes over $G^\dagger$.
\end{theorem}
\begin{proof}
Assume that $X$ is a Bayesian network on $G$. We can obviously rewrite the
probability distribution of $X$ in the form
$$
\pi(x) = \prod_{S\in G^\dagger_\CR}\prod_{s\in S} p_s(\pe x{\pa{s}}, \pe xs).
$$
Since every vertex in $S$ has its parents in $S$ or in $\bigcup_{T\in \pa{S}} T$,
this {\em a fortiori} takes the form
$$
\pi(x) = \prod_{S\in G^\dagger_\CR} p_S((\pe xT, T\in S^{-}), \pe xs ).
$$
So $\pe XS, S\in V_\CR$ is a Bayesian network. Moreover, 
$$
p_S((\pe xT, T\in S^{-}), \pe xs ) = \prod_{s\in S} p_s(\pe x{\pa{s}}, \pe xs)
$$
is a tree distribution with the required form of the individual
conditional distributions. 

Now assume that $X$ decomposes on $G^\dagger$. Then the conditional
distribution of $\pe XS$ given $\pe XT, T\in \pa{S}$ is Markov for the acyclic
undirected graph $G_S$, and can therefore be expressed as a tree
distribution consistent with the orientation of $G$. 
\end{proof}

\subsection{Markov equivalence}

While the previous discussion provides a rather simple description of
Bayesian networks in terms of chain graphs, it does not go all the way
in reducing the number of oriented edges in the definition of  a Bayesian
network. The issue is, in some way, addressed by the notion of Markov
equivalence, which is defined as follows.

\begin{definition}
\label{def:mar.eq}
Two directed acyclic graphs on the same set of vertexes $G = (V,E)$
and $\tilde G = (V, \tilde E)$ are Markov-equivalent
if any family of random variables that decomposes as a (positive) Bayesian network
over one of them also decomposes as a Bayesian network over the other.
\end{definition}

The notion of Markov equivalence is exactly described by
d-separation. This is stated in the following theorem, due to Geiger
and Pearl \citep{Geiger1990,geiger1990logic}, that we state without proof.
\begin{theorem}
\label{th:geig.p}
$G$ and $\tilde G$ are Markov equivalent if and only if, whenever two
vertexes are d-separated by a set in one of them, the same separation
is true with the other. 
\end{theorem}

This property can be expressed in a strikingly simple
condition. One says that a v-junction $(s,t,u)$ in a DAG is {\em unlinked}
if $s$ and $u$ are not neighbors.   
\begin{theorem}
\label{th:mar.eq}
$G$ and $\tilde G$ are Markov equivalent if and only if $G^\flat =
\tilde G^\flat$ and $G$ and $\tilde G$ have the same unlinked v-junctions.
\end{theorem}
\begin{proof}

\begin{enumerate}[label={\it Step \arabic*.}]
\item
We first show that
a given pair of vertexes in a DAG is unlinked if and only if it can
be d-separated by some set in the graph. Clearly, if they are linked,
they cannot be d-separated (which is the ``if'' part), so what really needs to be proved is that
unlinked vertexes can be d-separated. Let $s$ and $t$ be these vertexes
and let $U = \CA_{\{s,t\}}\setm \{s,t\}$. Then $U$ d-separates $s$ and
$t$ since any path between $s$ and $t$ in $(G_{\CA_{\{s,t\}\cup
U}})^\sharp = (G_{\CA_{\{s,t\}}})^\sharp$ must obviously pass in $U$.

\item  We now prove the only-if part of 
\cref{th:mar.eq} and
therefore assume that $G$ and $\tilde G$ are Markov equivalent, or, as
stated in  \cref{th:geig.p}, that
d-separation coincides in $G$ and $\tilde G$. We want to prove that
$G^\flat = \tilde G^\flat$ and unlinked v-junctions are the same.

\begin{enumerate}[label={\it Step 2.\arabic*.}]
\item The first statement is obvious from Step 1: d-separation determines
the existence of a link, so if d-separation coincides in the two
graphs, then the same holds for  links and $G^\flat = \tilde G^\flat$. 

\item So let us proceed to the second statement and
let $(s,t,u)$ be an unlinked
v-junction in $G$. We want to show that it is also a v-junction in
$\tilde G$ (obviously unlinked since links coincide). 

We will denote by $\tilde \CA_S$ the ancestors of some set $S\sub V$ in $\tilde G$ (while $\CA_S$
still denotes its ancestors in $G$). Let  $U =
\CA_{\{s,u\}}\setm \{s,u\}$.  Then,
as we have shown in Step 1, $U$ d-separates $s$ and $u$ in $G$,
so that, by assumption it also d-separates  them in $\tilde G$. 

We
know that $t\not \in U$, because it cannot be both a child and an
ancestor of $\{s, u\}$ in $G$ (this would induce a loop). The path
$(s,t,u)$ links $s$ and $u$ and does not pass in $U$, which is
only possible (since $U$ d-separates $s$ and $t$ in $\tilde G$) if it passes in
$V-\tilde \CA_U$ at a v-junction: so $(s,t,u)$ is a v-junction in
$\tilde G$, which is what we wanted to prove.
\end{enumerate}


\item We now consider the converse statement and assume that $G^\flat = \tilde G^\flat$
and unlinked v-junctions coincide. We want to show that d-separation
is the same in $G$ and $\tilde G$. So, we assume that $U$ d-separates $s$ and
$t$ in $G$, and we want to show that the same is true in $\tilde
G$. Thus, what we need to prove is:

{\em Claim 1.} Consider a path $\ga$ between $s$ and $t$ in $\tilde
G^\flat= G^\flat$. Then $\ga$ either  (D1)
passes in $U$ without a v-junction in $\tG$, or (D2) in $V\setm
\tilde\CA_U$ with a v-junction in $\tilde G$. 

We will prove Claim 1 using a series of lemmas. We say that $\ga$ has a three-point loop at
$u$ if $(v, u, w)$ are three consecutive points in $\ga$ such that
$v$ and $w$ are linked. So $(v, u, w, v)$ forms a loop in the
undirected graph.

\begin{lemma}
\label{lem:me.1}
If
$\gamma$ is a path between $s$ and $t$ that does not satisfy (D2) for $G$
and passes in $U$ without three-point loops, then $\ga$ satisfies
(D1) for $\tilde G$. 
\end{lemma}
The proof is easy: since $\ga$ does not satisfy (D2) in $G$,
it satisfies (D1) and passes in $U$ without a v-junction in
$G$. But this intersection cannot be a v-junction in $\tilde G$
since it would otherwise have to be linked and constitute a three-point
loop in $\ga$, which proves that (D1) is true for $\ga$ in $\tG$.

The next step is to remove the three-point loop condition in 
\cref{lem:me.1}. This will be done using the next two results.
\begin{lemma}
\label{lem:me.2}
Let $\ga$ be a path with a three-point loop at
$u\in U$ for $G$. Assume that $\ga\setm u$ (which is a valid path in
$G^\flat$) satisfies (D1) or (D2) in
$\tilde G$. Then $\ga$ satisfies (D1) or (D2) in $\tilde G$.
\end{lemma}
To prove the lemma, let $v$ and $w$ be the predecessor and
successor of $u$ in $\ga$. First assume that $\ga\setm u$ satisfies (D1) in
$\tilde G$. If
this does not happen at $v$ or at $w$, then this will apply also to
$\ga$ and we are done, so let us assume that $v\in U$ and that $(v', v,w)$ is not a
v-junction in $\tilde G$, where $v'$ is the predecessor of $v$.  If $(v',v,u)$ is
not a v-junction in $\tilde G$, then (D1) is true for $\ga$ in $\tilde
G$. If it is
a v-junction, then $(v,u,w)$ is not and (D1) is true too.

Assume now that (D2) is true for $\ga\setm u$ in $\tG$. Again, there
is no problem if (D2) occurs for some point other than $v$ or $w$, so
let us consider the case for which it happens at $v$. This means that 
$v\not\in \tilde \CA_U$ and $(v',v,w)$ is a v-junction. But, since
$u\in U$, the link between $u$ and $v$ must be from $u$ to $v$ in
$\tilde G$ so that
there is no v-junction at $u$ and (D1) is true in $\tilde G$. This proves 
\cref{lem:me.2}.

\begin{lemma}
\label{lem:me.3}
Let $\ga$ be a path with a three-point loop at
$u\in U$ for $G$. Assume that $\ga$ does not satisfy (D2) in $G$. Then
$\ga\setm u$ does not satisfy this property either.
\end{lemma}
Let us assume that $\ga\setm u$ satisfies (D2) and reach a
contradiction. Letting $(v,u,w)$ be the three-point loop, (D2) can only
happen in $\ga\setm u$ at $v$ or $w$, and let us assume that this
happens at $v$, so that, $v'$ being the predecessor of $v$, $(v',v,w)$
is a v-junction in $G$ with $v\not\in \CA_U$. Since $v\not\in \CA_U$,
the link between $u$ and $v$ in $G$ must be from $u$ to $v$, but this
implies that $(v',v,u)$ is a v-junction in $G$ with $v \not \in \CA_U$
which is a contradiction: this proves  \cref{lem:me.3}.


The previous three lemmas directly imply the next one.
\begin{lemma}
\label{lem:me.4}
If
$\gamma$ is a path between $s$ and $t$ that does not satisfy (D2) for
$G$, then $\ga$ satisfies
(D1) or (D2) for $\tilde G$. 
\end{lemma}
Indeed, if we start with $\ga$ that does not satisfy (D2) for $G$, 
 \cref{lem:me.3} allows us to progressively remove three-point loops from
$\ga$ until none remains with a final path that satisfies the
assumptions of  \cref{lem:me.1} and therefore satisfies (D1) in $\tilde G$,
and  \cref{lem:me.2} allows us to add the points that we have
removed in reverse order while always satisfying (D1) or (D2) in $\tG$.


We now partially relax the hypothesis that (D2) is not satisfied with the next lemma.
\begin{lemma}
\label{lem:me.5}
If
$\gamma$ is a path between $s$ and $t$ that does not pass in $V\setm
\CA_U$ at a linked v-junction for $G$, then $\ga$ satisfies
(D1) or (D2) for $\tilde G$. 
\end{lemma}
Assume that $\ga$  does not satisfy (D2) for
$\tilde G$ (otherwise the result is proved). By  \cref{lem:me.4},  $\ga$ must satisfy
(D2) for $G$. So, take an intersection of $\ga$ with $V\setm \CA_U$ that
occurs at a v-junction in $G$, that we will denote $(v, u, w)$. This is
still a v-junction in $\tilde G$ since we assume it to be
unlinked. Since (D2) is false in $\tilde G$, we must have
$u\in \tilde \CA_U$, and there is an oriented path, $\tau$, from
$u$ to $U$ in $\tilde G$. 

We can assume that $\tau$ has no v-junction in $G$. If a v-junction
exists in $\tau$,
then this v-junction must be linked (otherwise this would also be a
v-junction in $\tilde G$ and contradict the fact that $\tau$ is
consistently oriented in $\tilde G$), and this link must be oriented
from $u$ to $U$ in $\tilde G$ to avoid creating a loop in this
graph. This implies that we can bypass the v-junction while keeping a
consistently oriented path in $\tilde G$, and iterate this until
$\tau$ has no v-junction in $G$. But this implies that $\tau$ is
consistently oriented in $G$, necessarily from $U$ to $u$ since
$u \not\in \CA_U$.

Denote $\tau = (u_0=u, v_1, \ldots,
u_n\in U)$. We now prove by induction that each $(v, u_k, w)$ is an unlinked
v-junction. This is true when $k=0$, and let us assume that it is true
for $k-1$. Then $(u_{k}, u_{k-1}, v)$ is a v-junction in $G$ but not
in $\tilde G$: so it must be linked and there exists an edge between
$v$ and $u_k$. In $\tilde G$, this edge must be oriented from $v$
to $u_k$, since $(v, u_{k-1}, u_k, v)$ would form a loop
otherwise. For the same reason, there must be an edge in $\tilde G$
from $w$ to $u_k$ so that $(v, u_k, w)$ is an unlinked
v-junction. 

Since this is true for $k=n$, we can replace $u$ by $u_n$ in $\ga$
and still obtain a valid path. This can be done for all 
intersections of $\ga$ with $V\setm \CA_U$ that occur at
v-junctions. This finally yields a path (denote it
$\bar\ga$) which does not satisfy (D2) in $G$ anymore, and therefore
satisfies (D1) or (D2) in $\tG$: so  $\bar\ga$
must either pass in $U$ without a v-junction or in $V\setm \tilde A_U$
at a v-junction. None of the  nodes that were modified can satisfy any of these
conditions, since they were all in $U$ with a v-junction, so that the
result is true for the original $\ga$ also. This proves 
\cref{lem:me.5}.

So the only unsolved case is when $\ga$ is allowed to pass in $V\setm
\CA_U$ at linked v-junctions. We define an algorithm that removes them
as follows. Let $\ga_0=\ga$ and let $\ga_k$ be the path after step $k$ of
the algorithm. One passes from $\ga_k$ to $\ga_{k+1}$ as follows.
\begin{itemize}
\item If $\ga_k$ has no linked v-junctions in $V\setm \CA_U$ for $G$,
  stop.
\item Otherwise, pick such a v-junction and let $(v,u,w)$ be the
  three nodes involved in it.
\begin{enumerate}[label=(\roman*),labelindent=0.5cm, leftmargin=.5cm]
\item If $v\in U, v'\not\in U$ and $(v', v, u)$ is a v-junction in $\tG$, remove
  $v$ from $\ga_k$ to define $\ga_{k+1}$.
\item Otherwise,  if  $w\in U, w'\not\in U$ and $(u,w, w')$ is a v-junction in $\tG$, remove
  $w$ from $\ga_k$ to define $\ga_{k+1}$.
\item Otherwise, remove $u$ from $\ga_k$ to define $\ga_{k+1}$.
\end{enumerate}
\end{itemize}

None of the considered cases can disconnect the path. This is clear
for case (iii)  since $v$ and $w$ are linked. For case (i), note that,
in $G$, $(v',v, u)$ cannot be a v-junction since $(v,u,w)$ is
one. This implies that the v-junction in $\tG$ must be linked and that
$v'$ and $u$ are connected.

The algorithm will stop at some point with some $\ga_n$ that does not
have any linked v-junction in $V\setm \CA_U$ anymore, which implies
that (D1) or (D2) is true in $\tG$ for $\ga_n$. To prove that this
statement holds for $\ga$, it suffices to show that if (D1) or (D2) is true
in $\tG$ with
$\ga_{k+1}$, it must have been true with $\ga_k$ at each step of the
algorithm. So let's assume that $\ga_{k+1}$ satisfies (D1) or (D2) in $\tG$. 

First assume
that we passed from $\ga_k$ to $\ga_{k+1}$ via case (iii). Assume that
(D2) is true for $\ga_{k+1}$, with as usual the only interesting case
being when this occurs at $v$ or $w$. Assume it occurs at $v$ so that
$(v',v,w)$ is a v-junction and $v\not\in \tilde\CA_U$. If $(v',v,u)$
is a v-junction, then (D2) is true with $\ga_k$. Otherwise, there is an
edge from $v$ to $u$ in $\tG$ which also implies an edge from $w$ to
$u$ since $(v,u,w,v)$ would be a loop otherwise.  So $(v,u,w)$ is a
v-junction in $\tG$, and $u$ cannot be in $\tilde\CA_U$ since its
parent, $v$ would be in that set also. So (D2) is true in $\tG$. Now,
assume that 
(D1) is true at $v$, so that $(v',v,w)$ is not a v-junction and $v\in
U$. If $(v',v,u)$ is not a v-junction either, we are done, so assume
the contrary. If $v'\in U$, then we cannot have a v-junction at $v'$
and (D1) is true. But $v'\not\in U$ is not possible since this leads
to case (i).

Now assume that we passed from $\ga_{k}$ to $\ga_{k+1}$ via case
(i). Assume that (D1) is true for $\ga_k$: this cannot be at $v'$
since $v'\not\in U$, neither at $u$ since $u\not\in \CA_U$, so it will
also be true for $\ga_{k+1}$. The same statement holds with (D2) since
$(v',v,u)$ is a v-junction in $\tG$ with $v\in U$ which implies that
both $v'$ and $u$ are in $\tilde\CA_U$. Case (ii) is obviously
addressed similarly.
\end{enumerate}
With this,  the proof
of  \cref{th:mar.eq} is complete.
\end{proof}

 
\subsection{Probabilistic inference: Sum-prod algorithm}

We now discuss the issue of using the sum-prod algorithm to compute marginal probabilities,
$\myP(\pe Xs = \pe xs)$ for $s\in V$ when $X$ is a Bayesian network on $G =
(V,E)$. By definition, $\myP(X = x)$ can be written in the form
$$
\myP(X=x) = \prod_{C\in\CC} \phi_C(\pe xC)
$$
where $\CC$ contains all subsets $C_s := \defset{s} \cup \pa{s}$, $s\in
V$. Marginal probabilities can therefore be computed easily when the
factor graph associated to $\CC$ is acyclic, according to 
\cref{prop:fct.grph.cvg}.  However, because of the specific form of the $\phi_C$'s (they are
conditional probabilities), the sum-prod algorithm can be analyzed
in more detail, and provide correct results even when the factor graph
is not acyclic.

The general rules for the sum-prod algorithm are
\[
\left\{
\begin{aligned}
&m_{sC}(\pe xs) \leftarrow \prod_{\tilde C, s\in \tilde C,  \tilde C\neq C} m_{\tilde C s}(\pe xs)\\
&m_{Cs}(\pe xs) \leftarrow \sum_{\pe y{C}: \pe ys = \pe xs} \phi_C(\pe y{C}) \prod_{t\in C\setm \{s\}} m_{tC}(\pe yt)
\end{aligned}
\right.
\]

They take a particular form for Bayesian networks, using the fact that a vertex $s$ belongs to $C_s$, and to
all $C_t$ for $t\in \cl{s}$.
\begin{eqnarray*}
m_{sC_s}(\pe xs) &\leftarrow& \prod_{t\in \cl{s}} m_{C_ts}(\pe xs),\\
m_{sC_t}(\pe xs) &\leftarrow& m_{C_ss}(\pe xs) \prod_{u\in \cl{s}, u\neq t}
m_{C_us}(\pe xs), \text{ for } t\in \cl{s},\\
&&\\
m_{C_ss}(\pe xs) &\leftarrow& \sum_{\pe y{C_s}, \pe ys = \pe xs} p_s(\pe y{\pa{s}}, \pe xs) \prod_{t\in
\pa{s}} m_{tC_s}(\pe yt),\\
m_{C_ts}(\pe xs) &\leftarrow& \sum_{\pe y{C_t}, \pe ys = \pe xs} p_t(\pe xs\wedge \pe y{\pa{s}\setminus\{t\}}, \pe yt) m_{tC_t}(\pe yt) \prod_{u\in
\pa{s}, u\neq t} m_{uC_t}(\pe yu), 
\\ && \text{ for } t\in \cl{s}.
\end{eqnarray*}

These relations imply that, if $\pa{s}=\emp$ ($s$ is a root), then
$m_{C_ss} = p_s(\pe xs)$. Also, if $\cl{s} = \emp$ ($s$ is a leaf) then
$m_{sC_s} = 1$. The following proposition shows that many of the
messages become constant over time.
\begin{proposition}
\label{prop:bp.bay.1}
All upward messages, $m_{sC_s}$ and $m_{C_ts}$ with $t\in \cl{s}$ become
constant (independent from $\pe xs$) in finite time.
\end{proposition}
\begin{proof}
This can be shown recursively as follows. Assume that, for a given
$s$,  $m_{tC_t}$ is
constant for all $t\in \cl{s}$ (this is true if $s$ is a leaf). Then,
\begin{eqnarray*}
m_{C_ts}(\pe xs) &\leftarrow& \sum_{pe y{C_t}, \pe ys = \pe xs} 
p_t(\pe xs\wedge \pe y{\pa{s}\setminus\{t\}}, \pe yt) m_{tC_t}(\pe yt) \prod_{u\in
\pa{s}, u\neq t} m_{uC_t}(\pe yu), \\
&=& m_{tC_t} \sum_{\pe y{C_t}, \pe ys = \pe xs} p_t(\pe xs\wedge \pe y{\pa{s}\setminus\{t\}}, \pe yt) \prod_{u\in
\pa{s}, u\neq t} m_{uC_t}(\pe yu)\\
&=& m_{tC_t} \sum_{\pe y{C_t\setm \{t\}}, \pe ys = \pe xs} \prod_{u\in
\pa{s}, u\neq t} m_{uC_t}(\pe yu)\\
&=& m_{tC_t} \prod_{u\in \pa{s}, u\neq t} \sum_{\pe yu} m_{uC_t}(\pe yu)
\end{eqnarray*}
which is constant. Now
$$
m_{sC_s}(\pe xs) \leftarrow \prod_{t\in \cl{s}} m_{C_ts}(\pe xs)
$$
is also constant. This proves that all $m_{sC_s}$ progressively become
constant, and, as we have just seen, this implies the same property
for $m_{C_ts}$, $t\in \cl{s}$.
\end{proof}

This proposition implies that, if initialized with constant messages
(or after a finite time), the sum-prod algorithm iterates
\begin{eqnarray*}
m_{sC_s} &\leftarrow& \prod_{t\in \cl{s}} m_{C_ts}\\
m_{C_ss}(\pe xs) &\leftarrow& \sum_{\pe y{C_s}, \pe ys = \pe xs} p_s(\pe y_{\pa{s}}, \pe xs) \prod_{t\in
\pa{s}} m_{tC_s}(\pe yt)\\
m_{sC_t}(\pe xs)  &\leftarrow& m_{C_ss}(\pe xs) \prod_{u\in \cl{s}, u\neq t}
m_{C_us}, \  \ t\in \cl{s}\\
m_{C_ts} &\leftarrow& m_{sC_s} \prod_{u\in \pa{t}, u\neq s}
\sum_{\pe yu} m_{uC_s}(\pe yu), \  \ t\in \cl{s}.
\end{eqnarray*}
From this expression, we can conclude
\begin{proposition}
\label{prop:bp.bay.2}
If the previous algorithm is first initialized with upward messages,
$m_{sC_s} = m_{C_ts}$ all equal to 1, and if  downward messages are
computed top down from the roots to the leaves, the obtained
configuration of messages is invariant for the sum-prod algorithm.
\end{proposition}
\begin{proof}
If all upward messages are equal to 1, then clearly, the downward
messages sum to 1 once they are updated from roots to leaves, and this
implies that the upward messages will remain equal to 1 for the next
round.  The obtained configuration is invariant since the downward
messages are recursively uniquely defined by their value at the roots.
\end{proof}
The downward messages, under the
previous assumptions, satisfy $m_{sC_t}(\pe xs)  = m_{C_ss}(\pe xs)$ for all
$t\in \cl{s}$ and therefore
\begin{equation}
\label{eq:mes.bay}
m_{C_ss}(\pe xs) = \sum_{\pe y{C_s}, \pe ys = \pe xs} \pi(\pe y{\pa{s}}, \pe xs) \prod_{t\in
\pa{s}} m_{C_tt}(\pe yt).
\end{equation}
Note that the associated ``marginals'' inferred by the sum-prod algorithm are 
$$
\sig_s(\pe xs) = \prod_{C, s\in C} m_{Cs}(\pe xs) = m_{C_ss}(\pe xs) 
$$
since $m_{C_ts}(\pe xs) = 1$ when $t\in \cl{s}$.

Although the sum-prod algorithm initialized with unit messages
converges to a stable configuration if run top-down, the obtained
$\sig_s$'s do not necessarily provide the correct single site marginals. There is a
situation for which this is true, however, which is when the initial
directed graph is singly connected, as we will see below. Before this,
let us analyze the complexity resulting from an iterative computation of
the marginal probabilities, similar to what we have done with trees.

We define the
depth of a vertex in $G$ as follows.
\begin{definition}
\label{def:depth}
Let $G=(V,E)$ be a DAG. The depth of a vertex $s$ in $V$ is defined
recursively by
\begin{itemize}[labelindent=0.5cm]
\item[-] $\mathrm{depth}(s) = 0$ if $s$ has no parent.
\item[-] $\mathrm{depth}(s) = 1 + \max\big(\mathrm{depth}(t), t\in
\pa{s}\big)$ otherwise.
\end{itemize}
\end{definition}

The recursive computation of marginal distributions is made possible
(although not always feasible)
with the following remark.
\begin{lemma}
\label{lem:c.i.dag}
Let $X$ be a Bayesian network on the DAG $G=(V,E)$, and $S\sub V$,
such that all elements in $S$ have the same depth. Let
$\pa{S}$ be the set of parents of elements in $S$, and $T =
\mathrm{depth}^-(S)$ the set of vertexes in $V$ with depth strictly smaller
than the depth of $S$. Then 
$\cind{\pe XS}{\pe X{T\setm \pa{S}}}{\pe X{\pa{S}}}$ and the variables $\pe Xs, s\in S$
are conditionally independent given $\pe X{\pa{S}}$. 
\end{lemma}
\begin{proof}
It suffices to show that vertexes in $S$ are separated from  $T\setm
\pa{S}$ and from other elements of $S$ by $\pa{S}$ for the
graph $(G_{S\cup T})^\sharp$. Any path starting at $s\in S$ must either
pass by a parent of $s$ (which is what we want), or by one of its children, or by another
vertex that shares a child with $s$ in $G_{S\cup T}$. But $s$ cannot have any
child in $G_{S\cup T}$, since this child cannot have a smaller depth
than $s$,
and it cannot be in $S$ either since all elements in $S$ have the same
depth.
\end{proof}
This lemma allows us to work recursively as follows. Assume that we can
compute marginal distributions over sets $S$ with maximal depth no
larger than $d$. Take a set $S$ of maximal depth $d+1$, and let $S_0$ be the
set of elements of depth $d+1$ in $S$. Then, letting
$T=\mathrm{depth}^-(S) = \mathrm{depth}^-(S_0)$, and $S_1 = S\setm S_0$,
\begin{eqnarray}
\nonumber
\myP(\pe XS = \pe xS) &=& \sum_{\pe y{T\setm S_1}} \myP(\pe X{S_0} = \pe x{S_0} \mid \pe XT = \pe y{T\setm S_1}\wedge
\pe x{S_1}) P(\pe X{T\cup S_1} =  \pe y{T\setm S_1}\wedge \pe x{S_1})\\
\label{eq:rec.inf}
&=& \sum_{\pe y{\pa{S}\setm S_1}} \prod_{s\in S_0} p_s(\pe {(y\wedge x)}{\pa{s}}\wedge
\pe x{S_1}, \pe xs) \myP(\pe X{\pa{S_0}\cup S_1} = \pe  y{\pa{S_0}\setm S_1}\wedge \pe x{S_1})
\end{eqnarray}
Since $\pa{S}\cup S_1$ has maximal depth strictly smaller than the
maximal depth of $S$, this indeed provides a recursive formula for the
computation of marginal over subsets of $V$ with increasing maximal
depths. However, because one needs to add parents to the
considered set when reducing the depth, one may end up having to
compute marginals over very large sets, which becomes intractable
without further assumptions.

A way to reduce the complexity is to assume that the graph $G$ is
singly connected, as defined below.

\begin{definition}
\label{def:sin.com}
A DAG $G$ is singly connected if there exists at most one path in $G$ that connects
any two vertexes. 
\end{definition}
Such a property is true for a tree, but also
holds for some networks with multiple parents. We have the
following nice property in this case.

\begin{proposition}
\label{prop:sin.con}
Let $G$ be a singly connected DAG and $X$ a Bayesian network on
$G$. If $s$ is a vertex in $G$, the variables $(\pe Xt, t\in \pa{s})$ are
mutually independent.
\end{proposition}
\begin{proof}
We have, using  \cref{prop:rest},
$$
\myP(\pe X{\pa{s}} = \pe x{\pa{s}})  = \sum_{\pe y{\CA_{\pa{s}}}, \pe y{\pa{s}} = \pe x{\pa{s}}}
\prod_{u\in \CA_{\pa{s}}} p_u(\pe y{\pa{u}}, \pe yu).
$$
Because the graph is singly connected, two parents of $s$ cannot have a
common ancestor (since there would then be two paths from this
ancestor to 
$S$). So $\CA_{\pa{s}}$ is the disjoint union of the $\CA_{t}$'s for $t\in
\pa{s}$ and we can write
\begin{eqnarray*}
\myP(\pe X{\pa{s}} = \pe x{\pa{s}})  &=& \sum_{\pe y{\CA_{\pa{s}}}, \pe y{\pa{s}} = \pe x{\pa{s}}}
 \prod_{t\in \pa{s}}  \prod_{u\in\CA_t} p_u(\pe y{\pa{u}}, \pe yu)\\
&=& \prod_{t\in \pa{s}} \sum_{\pe y{\CA_{t}}, \pe y{t} = \pe x{t}}
   \prod_{u\in\CA_t} p_u(\pe y{\pa{u}}, \pe yu)\\
&=& \prod_{t\in \pa{s}} \myP(\pe Xt = \pe xt)
\end{eqnarray*}
This proves the lemma.
\end{proof}

 \Cref{eq:rec.inf} can be simplified under the assumption of
a singly connected graph, at least for the computation of single
vertex marginals; we have, if $s\in V$ and $G$ is singly connected
\begin{equation}
\label{eq:sing.con}
\myP(\pe Xs = \pe xs) = \sum_{\pe y{\pa{s}}}  p_s(\pe y{\pa{s}}, \pe xs) \prod_{t\in \pa{s}}
\myP(\pe X{t} = \pe yt).
\end{equation}
This is now recursive in single vertex marginal probabilities. It
moreover coincides with the recursive equation that defines the
messages $m_{C_ss}$ in \cref{eq:mes.bay}, which shows that the
sum-prod algorithm provides the correct answer in this case.

\subsection{Conditional probabilities and interventions}
One of the main interests of graphical models is to provide an
ability to infer the behavior of hidden variables of interest given
other, observed, variables. When dealing with oriented graphs the
way this should be analyzed is, however, ambiguous. 

Let's consider an example, provided by the graph in  \cref{fig:snow}.
\begin{figure}[h]
\centering
		\begin{tikzpicture}[mynode/.style={ellipse, draw=black!60, fill=gray!5, ultra thick, , minimum size=10mm},]
				\node[mynode](H1){No school};
				\node[mynode](H2)[above left =1.5cm and -.5cm of H1]{Bad weather};
				\node[mynode](H3)[above right=1.5cm and -.5cm of H1]{Broken HVAC};

		\draw[black, ultra thick,->] (H2.south) -- (H1.north);
		\draw[black, ultra thick,->] (H3.south) -- (H1.north);
		\end{tikzpicture}
%\includegraphics{graph1.pdf}
\caption{\label{fig:snow} Example of causal graph.}
\end{figure}
The Bayesian network interpretation of this graph is that both
events (which may be true or false) ``Bad weather'' and ``Broken HVAC'' happen first, and that they are
independent. Then, given their observation, the ``No school'' event may
occur, probably more likely if the weather is bad or the HVAC is broken or snow, and even more likely if both happened at the same time.

Now consider the following passive observation: you wake up, you
haven't checked the weather yet or the news yet, and someone tells you that there is no
school today. Then you may infer that there is more chances than usual
for bad weather or the HVAC broken at school. Conditionally to this information, these two events become
correlated, even if they were initially independent. So, even if the
``No school'' event is considered as a probabilistic consequence of its
parents, observing it influences our knowledge on them.

Now, here is an intervention, or manipulation: the school superintendent has declared
that he has given enough snow days for the year and declared that there
would be school today whatever happens. So you know that the
``no-school'' event will not happen. Does it change the risk of bad weather of broken HVAC? Obviously not: an intervention on a node does not affect the
distribution of the parents. 

Manipulation and passive observation are two very different ways of
affecting unobserved variables in Bayesian networks. Both of them may
be relevant in applications. Of the two, the simplest to analyze is
intervention, since it merely consists in clamping one of the
variables while letting the rest of the network dynamics
unchanged. This leads to the following formal definition of
manipulation.
\begin{definition}
\label{def:manip}
Let $G = (V,E)$ be a directed acyclic graph and $X$ a Bayesian network
on $G$. Let $S$ be a subset of $G$
and $\pe xS\in F_S$ a given configuration on $S$. Then the manipulated distribution
of $X$ with fixed values $\pe xS$ on $S$ is the Bayesian network on the
restricted graph $G_S$, with the same conditional probabilities, using
the value $\pe xs$ every time a vertex $s\in S$ is a parent of $t\in
V\setm S$ in $G$.
\end{definition}
So, if the distribution of $X$ is given by \cref{eq:bn.2}, then its
distribution after manipulation on $S$ is 
$$
\tilde \pi(\pe y{V\setm S}) = \prod_{t\in V\setm S} p_t(\pe y{\pa{t}}, \pe yt)
$$
where $\pa{t}$ is the set of parents of $t$ in $G$, and $\pe ys=\pe xs$
whenever $s\in \pa{t}\cap S$.


The distribution of a Bayesian network $X$ after passive observation
$\pe XS = \pe xS$ is
not so easily described. It is obviously the conditional distribution
$P(\pe X{V\setm S} = \pe y{V\setm S}\mid \pe XS = \pe xS)$ and therefore requires using
the conditional dependency structure, involving the moral
graph and/or d-separation.

Let us discuss this first in the simpler case of trees, for which the
moral graph is the undirected acyclic graph underlying the tree, and
d-separation is simple separation on this acyclic graph. We can then
use  \cref{prop:cond.dist.grph} to understand the new
structure after conditioning: it is a $G^\flat_{V\setm S}$-Markov
random field, and, for $t\in V\setm S$, the conditional distribution
of $\pe Xt = \pe yt$ given its neighbors is the same as before, using the
value $\pe xs$ when $s\in S$. But note that when doing this (passing to $G^\flat$), we broke the
causality relation between the variables. We can however always go
back to a tree (or forest, since connectedness may have been broken) with
the same edge orientation as they initially were, but this requires
reconstituting the edge joint probabilities from the new acyclic
graph, and therefore using (acyclic) belief propagation.

With general Bayesian networks, we know that the moral graph can be
loopy and therefore a source of  difficulties. The following
proposition states that the damage is circumscribed to the ancestors of
$S$. 

\begin{proposition}
\label{prop:inf.cond}
Let $G=(V,E)$ be a directed acyclic graph, $X$ a Bayesian network on
$G$, $S\sub V$ and $\pe x{\CA_S}\in \CF(\CA_S)$. Then the conditional
distribution of $\pe X{\CA_S^c}$ given by $\pe X{\CA_S} = \pe x{\CA_S}$
coincides with the manipulated distribution in 
\cref{def:manip}.
\end{proposition}
\begin{proof} 
The conditional distribution is proportional to 
$$
\prod_{s\in V} p(\pe y{\pa{s}}, \pe ys)
$$
with $\pe yt = \pe xt$ if $t\in \CA_S$. Since $s\in\CA_S$ implies $\pa{s}\sub
\CA_S$, all terms with $s\in \CA_S$ are constant in the sum and can be
factored out after normalization. So the conditional distribution is
proportional to 
$$
\prod_{s\in \CA_S^c} p(\pe y{\pa{s}}, \pe ys)
$$
with $\pe yt = \pe xt$ if $t\in \CA_S$. But we know that such products sum
to 1, so that the conditional distribution is equal to this expression
and therefore provides a Bayesian network on $G_{\CA_S^c}$.
\end{proof}

\section{Structural equation models}
\label{sec:sem}

Structural equation models (SEM's) provides an alternative (and essentially equivalent) formulation of Bayesian networks, which may be more convenient to use, especially when dealing with variables taking values in general state spaces.

Let $G=(V, E)$ be a directed acyclic graph.  SEMs are associated to families of functions $\Phi_s: \CF(\pa{s}) \times \CB_s\to F_s$ and random variables $\bfxi_s: \Omega \to \CB_s$ (where $\CB_s$ is some measurable set), for $s\in V$. The random field $X: \Omega \to \CF(V)$ associated to the SEM satisfies the equations
\begin{equation}
\label{eq:sem}
X_s =\pe \Phi s(\pe X{s-}, \pe \bfxi s).
\end{equation}
Because of the DAG structure, these equations uniquely define $X$ once $\bfxi$ is specified. As a consequence, there exists a function $\boldsymbol\Psi$ such that $X = \boldsymbol \Psi(\bfxi)$. 


%The model is generally parametrized by adding parameters to the functions $\pe \Phi s$, and we will use the notation $\pe {\Phi_\theta}s$, where $\theta$ parametrizes all structural equations.  
The model is therefore fully specified by the functions $\pe \Phi s$ and the probability distributions of the variables $\pe \bfxi s$. We will assume that they have a density, denoted $\pe g s, s\in V$, with respect to some measure $\mu_s$ on $\CB_s$. 
%We do not include parameters in these densities, to simplify the notation. 
They are typically chosen as uniform distributions on $\CB_s$ (continuous and compact, or discrete) or as standard Gaussian when $\CB_s = \mR^{d_s}$ for some $d_s$. One also generally assumes that the variables $(\pe \bfxi s, s\in V)$ are jointly independent, and we make this assumption below.

%If $\CB_s = \mR^{d_s}$ and $\mu_s$ is Lebesgue's measure, the 

Let $V_k$, $k\geq 0$, be the set of vertexes in $V$ with depth $k$ (c.f. \cref{def:depth}) and $V_{<k} = V_0 \cup \cdots\cup V_{k-1}$. Then (using the independence of $(\pe\bfxi s, s\in V)$, for $s\in V_k$,  the conditional distribution of $\pe Xs$ given $\pe X{V_{<k}} = \pe x{V_{<k}}$ is the distribution  of $\pe\Phi s(\pe x{s-}, \pe\bfxi s)$. Formally this is given by 
\[
\pe\Phi s(\pe x{s-}, \cdot)_{\sharp} (\pe gs\mu_s),
\]
the pushforward of the distribution of $\pe \bfxi s$ by $\pe\Phi s(\pe x{s-}, \cdot)$.

More concretely, assume that $\bfxi_s$ follows a uniform distribution on $\CB_s = [0, 1]^h$ for some $h$, and assume that $F_s$ is finite for all $s$. Then,
\[
P(\pe Xs = \pe xs \mid \pe X{V_{<k}} = \pe x{V_{<k}}) = \mathrm{Volume}(U_s(  \pe x{\pa{s}}, \pe xs)) \defeq p_s( \pe x{\pa{s}}, \pe xs)
\]
where 
\[
U_s(\pe x{\pa{s}}, \pe xs) = \left\{\xi\in [0,1]^h :  \pe\Phi s(\pe x{s-}, \xi) = \pe xs\right\}.
\] 
Since variables $\pe Xs$, $s\in V_k$ are conditionally independent given $\pe X {V<k}$, we find that $X$ decomposes as a Bayesian network over $G$, 
\[
P(X=x) = \prod_{s\in V}  p_s( \pe x{\pa{s}}, \pe xs).
\]
Similarly, if $F_s = \CB_s = \mR^{d_s}$, $\pe \xi s\sim\CN(0, \Id[d_s])$, and $\pe \bfxi s \mapsto \pe{\Phi_\theta}s(\pe x{\pa{s}}, \pe \xi s)$ is invertible, with $C^1$ inverse $\pe xs \mapsto \pe{\Psi_\theta} s(\pe x{\pa{s}}, \pe xs)$, then $X$ is a Bayesian network, with continuous variables, and, using the change of variable formula, the conditional distribution of $\pe Xs$ given $\pe X{\pa{s}} = \pe x{s-}$ has  p.d.f.
\[
 p_s( \pe x{\pa{s}}, \pe xs) = \frac{1}{(2\pi)^{d_s/2}} \exp\left(-\frac12 |x_s - 
\pe{\Psi_\theta}s(\pe x{\pa{s}}, \pe xs)|^2\right) \left|\det (\partial_{\pe xs}\pe{\Psi_\theta}s(\pe x{\pa{s}}, \pe xs))\right|.
\]
A simple and commonly used special case for this example are linear SEMs, with
\[
\pe Xs = a_s + b_s^T\pe Xs + \sigma_s \pe\xi s.
\] 
In this case, the inverse mapping is immediate and the Jacobian determinant in the change of variables is $1/\sigma_s^{d_s}$. 
%Since they are Bayesian networks, maximum likelihood estimators  for SEMs can be learned, in principle, using the methods discussed in \cref{chap:bayes.net}. %However, alternative learning paradigms have been introduced, as discussed in the next sections. 


